% GENERATED FILE. Edit canonical lessons, then run npm run generate:books.
% canonical-source-hash: fnv1a64:a8f05304ad177b35
% canonical-lessons: LA-C245-fur, LA-C245-latro, LA-C245-consul, LA-C245-imperator, LA-C245-regnum

\chapter{Thief, Robber, Consul, General, Kingdom}
\label{ch:la-thief-robber-consul-general-kingdom}
\hlchaptermodality{245}

\begin{chapteropening}
\textbf{By the end of this chapter:} \emph{I can say a thief, a robber, a consul, a general and a kingdom.}

Complete the last lesson of chapter 245: I can say a thief, a robber, a consul, a general and a kingdom.
\end{chapteropening}

\section[fūr, fūris]{fūr, fūris — a thief}

\label{lesson:LA-C245-fur}



\begin{quote}
Before the new one: say the Latin for the public baths, then the Latin for an amphitheatre.
\end{quote}

\subsection*{You'll want to know: fūr, fūris}
\textbf{fūr, fūris} — \textquotedblleft{}a thief\textquotedblright{}.

\begin{cousinweb}
\textbf{fūr} is inherited, usually traced to Proto-Indo-European \emph{*bhōr}, \textquotedblleft{}thief\textquotedblright{} (EtymDB 2.1), from the root of \emph{ferō}, \textquotedblleft{}to carry\textquotedblright{}; Greek \emph{phōr}, \textquotedblleft{}thief\textquotedblright{}, is a cousin (Lewis \& Short).
\end{cousinweb}

\begin{grammarlens}[title={What you've built}]
One more word for this chapter.
\end{grammarlens}

\subsection*{Your turn}
\begin{itemize}
\raggedright
  \item \emph{Say it:} \emph{fūr, fūris}
  \item \emph{Say it:} \emph{fūr, fūris}, once more
  \item \emph{Say it:} say \emph{fūr, fūris}
  \item \emph{Recall:} say the Latin for the public baths, then the Latin for an amphitheatre, then say \emph{fūr, fūris} again
\end{itemize}

\begin{tcolorbox}[breakable,colback=teal!4,colframe=teal!35!black,title={Before you move on}]
What does fūr, fūris mean? (\textquotedblleft{}A thief\textquotedblright{}.) Say it once more.
\end{tcolorbox}
\section[latrō, latrōnis]{latrō, latrōnis — a robber}

\label{lesson:LA-C245-latro}



\begin{quote}
Before the new one: say the Latin for an amphitheatre, then the Latin for a thief.
\end{quote}

\subsection*{You'll want to know: latrō, latrōnis}
\textbf{latrō, latrōnis} — \textquotedblleft{}a robber\textquotedblright{}.

\begin{cousinweb}
\textbf{latrō} was probably borrowed from Greek: \emph{latron}, \textquotedblleft{}pay, hire\textquotedblright{} (EtymDB 2.1), and \emph{latris}, \textquotedblleft{}a hired servant\textquotedblright{} (Lewis \& Short). A \emph{latrō} was first a hired soldier, then a robber.
\end{cousinweb}

\begin{grammarlens}[title={What you've built}]
One more word for this chapter.
\end{grammarlens}

\subsection*{Your turn}
\begin{itemize}
\raggedright
  \item \emph{Say it:} \emph{latrō, latrōnis}
  \item \emph{Say it:} \emph{latrō, latrōnis}, once more
  \item \emph{Say it:} say \emph{latrō, latrōnis}
  \item \emph{Recall:} say the Latin for an amphitheatre, then the Latin for a thief, then say \emph{latrō, latrōnis} again
\end{itemize}

\begin{tcolorbox}[breakable,colback=teal!4,colframe=teal!35!black,title={Before you move on}]
What does latrō, latrōnis mean? (\textquotedblleft{}A robber\textquotedblright{}.) Say it once more.
\end{tcolorbox}
\section[cōnsul, cōnsulis]{cōnsul, cōnsulis — a consul}

\label{lesson:LA-C245-consul}



\begin{quote}
Before the new one: say the Latin for a thief, then the Latin for a robber.
\end{quote}

\subsection*{You'll want to know: cōnsul, cōnsulis}
\textbf{cōnsul, cōnsulis} — \textquotedblleft{}a consul\textquotedblright{}.

\begin{cousinweb}
\textbf{cōnsul} belongs with \emph{cōnsulō}, \textquotedblleft{}to consult, to deliberate\textquotedblright{} (EtymDB 2.1); Old Latin wrote it \emph{consol}. Where the root comes from is uncertain.
\end{cousinweb}

\begin{grammarlens}[title={What you've built}]
One more word for this chapter.
\end{grammarlens}

\subsection*{Your turn}
\begin{itemize}
\raggedright
  \item \emph{Say it:} \emph{cōnsul, cōnsulis}
  \item \emph{Say it:} \emph{cōnsul, cōnsulis}, once more
  \item \emph{Say it:} say \emph{cōnsul, cōnsulis}
  \item \emph{Recall:} say the Latin for a thief, then the Latin for a robber, then say \emph{cōnsul, cōnsulis} again
\end{itemize}

\begin{tcolorbox}[breakable,colback=teal!4,colframe=teal!35!black,title={Before you move on}]
What does cōnsul, cōnsulis mean? (\textquotedblleft{}A consul\textquotedblright{}.) Say it once more.
\end{tcolorbox}
\section[imperātor, imperātōris]{imperātor, imperātōris — a general, an emperor}

\label{lesson:LA-C245-imperator}



\begin{quote}
Before the new one: say the Latin for a robber, then the Latin for a consul.
\end{quote}

\subsection*{You'll want to know: imperātor, imperātōris}
\textbf{imperātor, imperātōris} — \textquotedblleft{}a general, an emperor\textquotedblright{}.

\begin{cousinweb}
\textbf{imperātor} is built on \emph{imperō}, \textquotedblleft{}to command\textquotedblright{} (EtymDB 2.1), itself \emph{in}, \textquotedblleft{}in, upon\textquotedblright{}, with \emph{parō}, \textquotedblleft{}to prepare, to order\textquotedblright{}, with \emph{-tor}, \textquotedblleft{}one who does it\textquotedblright{}. English \emph{emperor} comes from it through French.
\end{cousinweb}

\begin{grammarlens}[title={What you've built}]
One more word for this chapter.
\end{grammarlens}

\subsection*{Your turn}
\begin{itemize}
\raggedright
  \item \emph{Say it:} \emph{imperātor, imperātōris}
  \item \emph{Say it:} \emph{imperātor, imperātōris}, once more
  \item \emph{Say it:} say \emph{imperātor, imperātōris}
  \item \emph{Recall:} say the Latin for a robber, then the Latin for a consul, then say \emph{imperātor, imperātōris} again
\end{itemize}

\begin{tcolorbox}[breakable,colback=teal!4,colframe=teal!35!black,title={Before you move on}]
What does imperātor, imperātōris mean? (\textquotedblleft{}A general, an emperor\textquotedblright{}.) Say it once more.
\end{tcolorbox}
\section[rēgnum, rēgnī]{rēgnum, rēgnī — a kingdom}

\label{lesson:LA-C245-regnum}



\begin{quote}
Before the new one: say the Latin for a consul, then the Latin for a general.
\end{quote}

\subsection*{You'll want to know: rēgnum, rēgnī}
\textbf{rēgnum, rēgnī} — \textquotedblleft{}a kingdom\textquotedblright{}.

\begin{cousinweb}
\textbf{rēgnum} is built on \emph{rēx}, \emph{rēgis}, \textquotedblleft{}a king\textquotedblright{} (Lewis \& Short): what a king rules. English \emph{reign} comes from it.
\end{cousinweb}

\begin{grammarlens}[title={What you've built}]
One more word for this chapter.
\end{grammarlens}

\subsection*{Your turn}
\begin{itemize}
\raggedright
  \item \emph{Say it:} \emph{rēgnum, rēgnī}
  \item \emph{Say it:} \emph{rēgnum, rēgnī}, once more
  \item \emph{Say it:} say \emph{rēgnum, rēgnī}
  \item \emph{Recall:} say the Latin for a consul, then the Latin for a general, then say \emph{rēgnum, rēgnī} again
\end{itemize}

\begin{tcolorbox}[breakable,colback=teal!4,colframe=teal!35!black,title={Before you move on}]
What does rēgnum, rēgnī mean? (\textquotedblleft{}A kingdom\textquotedblright{}.) Say it once more.
\end{tcolorbox}
